\documentclass[11pt,a4paper]{article}

\usepackage[provide=*,spanish]{babel}
\usepackage[utf8]{inputenc}
\usepackage[T1]{fontenc}
\usepackage{lmodern}
\usepackage[margin=2.35cm]{geometry}
\usepackage{amsmath}
\usepackage{booktabs}
\usepackage{longtable}
\usepackage{tabularx}
\usepackage{array}
\usepackage{enumitem}
\usepackage{listings}
\usepackage{xcolor}
\usepackage{microtype}
\usepackage{parskip}
\usepackage{titlesec}
\usepackage{hyperref}
\usepackage{graphicx}
\usepackage{caption}
\usepackage{subcaption}
\usepackage{float}

\definecolor{azulWCD}{RGB}{17,67,112}
\definecolor{azulClaro}{RGB}{232,241,249}
\definecolor{grisCodigo}{RGB}{247,247,247}
\definecolor{verdeCodigo}{RGB}{25,110,55}
\definecolor{rojoCodigo}{RGB}{145,35,35}
\definecolor{tituloColor}{RGB}{0,55,110}
\definecolor{reglaColor}{RGB}{0,90,160}

\hypersetup{
  colorlinks=true,
  linkcolor=azulWCD,
  urlcolor=azulWCD,
  citecolor=azulWCD,
  pdftitle={Dos campanas de simulacion: con y sin tratamiento termico S(alpha,beta)},
  pdfauthor={Jaime A. Betancourt}
}

\lstdefinestyle{pystyle}{
  language=Python,
  backgroundcolor=\color{grisCodigo},
  basicstyle=\ttfamily\footnotesize,
  keywordstyle=\color{azulWCD}\bfseries,
  commentstyle=\color{verdeCodigo}\itshape,
  stringstyle=\color{rojoCodigo},
  breaklines=true,
  breakatwhitespace=false,
  frame=single,
  rulecolor=\color{gray!40},
  tabsize=2,
  showstringspaces=false,
  columns=fullflexible,
  keepspaces=true
}
\lstset{style=pystyle}

\titleformat{\section}{\normalfont\Large\bfseries\color{azulWCD}}
  {\thesection.}{0.6em}{}
\titleformat{\subsection}{\normalfont\large\bfseries}
  {\thesubsection.}{0.5em}{}
\setlist{itemsep=0.25em,topsep=0.4em}
\renewcommand{\arraystretch}{1.18}
\setlength{\emergencystretch}{2em}

\newcommand{\codigo}[1]{\nolinkurl{#1}}
\newcommand{\archivo}[1]{\path{#1}}
\newcommand{\var}[1]{\texttt{#1}}

\begin{document}
\pagestyle{empty}

%------------------------------------------------
% PORTADA
%------------------------------------------------

\begin{titlepage}
\begin{center}

\vspace*{0.5cm}
\textcolor{reglaColor}{\rule{\linewidth}{1.2pt}}

\vspace{1.3cm}
{\Large\scshape Universidad Industrial de Santander}\\[0.2cm]
{\large Escuela de Física --- Facultad de Ciencias}\\[0.2cm]
{\large Doctorado en Física}

\vspace{2.2cm}
{\Huge\bfseries\color{tituloColor} Dos campañas de simulación:\\[0.2cm]
con y sin tratamiento térmico S($\alpha,\beta$)}

\vspace{1.8cm}
{\large\textbf{Introducción al repositorio de análisis}}

\vspace{0.6cm}
\begin{minipage}{0.82\linewidth}
\centering
\textit{``Léase antes de comparar cifras entre notebooks o informes de este repositorio''}
\end{minipage}

\vfill

{\large\textbf{Autor}}\\[0.15cm]
Jaime A. Betancourt\\[0.1cm]
Doctorado en Física, Universidad Industrial de Santander

\vspace{1.2cm}
\textcolor{reglaColor}{\rule{0.5\linewidth}{0.6pt}}

\vspace{1cm}
18 de septiembre de 2026

\vspace{1cm}
\textcolor{reglaColor}{\rule{\linewidth}{1.2pt}}

\end{center}
\end{titlepage}

\pagestyle{plain}

\begin{abstract}
Este documento introduce, para un lector sin conocimiento previo del desarrollo de la tesis, la existencia de \textbf{dos campañas de simulación Geant4/MEIGA distintas} del detector Cherenkov de agua (WCD) con agua pura y soluciones de NaCl (2.5\,\%, 5\,\%, 10\,\% en masa), documentadas en paralelo en este repositorio. Ambas comparten geometría, materiales de estructura y framework, pero difieren en si el transporte de neutrones subtérmicos incluye o no el tratamiento térmico $S(\alpha,\beta)$ para el hidrógeno ligado en moléculas de agua --y eso hace que sus resultados, en particular la fracción de captura, no coincidan a energías comparables. Cualquier comparación de cifras entre un notebook/informe y otro debe tener esta distinción presente. Este texto corresponde al mismo contenido publicado como introducción en \archivo{analysis/README.md} del repositorio.
\end{abstract}

\tableofcontents

\section{Qué documenta este repositorio}

Este repositorio contiene los programas y cuadernos utilizados para leer y depurar las salidas de simulación, construir histogramas con intervalos lineales o logarítmicos, normalizar distribuciones por neutrón incidente y ancho de intervalo, calcular eficiencias e incertidumbres estadísticas, comparar agua pura con soluciones de NaCl, y producir las tablas y figuras de la tesis.

\section{Dos campañas de simulación: léase antes de comparar resultados}
\label{sec:campanas}

Los notebooks y los informes técnicos de este repositorio documentan \textbf{dos campañas de simulación Geant4/MEIGA distintas} del mismo WCD. Comparten geometría, materiales de estructura y framework, pero \textbf{difieren en si el transporte de neutrones subtérmicos incluye o no el tratamiento térmico $S(\alpha,\beta)$} para el hidrógeno ligado en moléculas de agua --y eso hace que sus resultados, en particular la fracción de captura, no coincidan a energías comparables. Quien llegue a este repositorio sin haber seguido el desarrollo de la tesis debe tener esto presente antes de cruzar cifras entre un notebook/informe y otro:

\begin{table}[H]
\centering
\small
\begin{tabularx}{\linewidth}{@{}l X X@{}}
\toprule
 & \textbf{Campaña 1 --- ``primeras simulaciones''} & \textbf{Campaña 2 --- exploratoria} \\
\midrule
Tratamiento térmico $S(\alpha,\beta)$ & \textbf{No} implementado & \textbf{Sí} implementado \\
\addlinespace
Energías de inyección & 1\,meV, 10\,meV, 25\,meV, 100\,meV, 10\,keV (5 corridas $\times$ 4 medios = 20 corridas) & 25\,meV únicamente (1 corrida $\times$ 4 medios), por ahora \\
\addlinespace
Clasificación captura/reflexión/transmisión & Por geometría (posición del neutrón); el archivo de trazabilidad no registra material -- ver nota sobre volumen del tanque más abajo & Por material (el archivo de trazabilidad sí lo registra) \\
\addlinespace
Rol en la tesis & \textbf{Es la campaña sobre la que se apoya el análisis completo de la tesis} & Corrida de verificación/comparación, aún parcial \\
\addlinespace
Informe técnico & \codigo{campana-1\_sin-S-alpha-beta/Informe-tecnico} & \codigo{campana-2\_con-S-alpha-beta\_25meV/Informe-tecnico} \\
\addlinespace
Notebooks & Viven junto a los datos crudos, fuera de este repositorio (no replicados aquí por su tamaño) & \codigo{analysis/notebooks/campana-2\_con-S-alpha-beta\_25meV/} \\
\addlinespace
Carpeta de datos crudos (fuera del repo) & \codigo{Nuevas-simulaciones-2026/Primeras-simulaciones/} (incluye la subcarpeta mal etiquetada \codigo{1keV/}, que en realidad es 10\,keV -- ver nota abajo) & \codigo{Nuevas-simulaciones-2026/25meV/} (nombre correcto -- 25\,meV, verificado directamente en \codigo{neutrones-incidentes.tsv} de las 4 concentraciones) \\
\bottomrule
\end{tabularx}
\caption{Comparación entre las dos campañas de simulación documentadas en este repositorio.}
\label{tab:campanas}
\end{table}

\subsection{Por qué importa esta distinción}

Ambas campañas comparten el punto de \textbf{25\,meV} --elegido en la Campaña 2 precisamente para
poner a prueba el régimen térmico, en torno a la energía térmica del medio, donde $S(\alpha,\beta)$
más debería notarse. $S(\alpha,\beta)$ modifica la sección eficaz de dispersión elástica del
hidrógeno ligado justo en ese rango de energías, así que \textbf{25\,meV es el punto natural de
comparación directa entre las dos campañas} --no debe asumirse que el balance de
captura/reflexión/transmisión ni la letargía de una reproducen los de la otra ahí. Los primeros
resultados ya muestran una fracción de captura distinta entre campañas a esa misma energía.
Cuantificar sistemáticamente esa diferencia a partir de este punto --y extender la Campaña 2 a las
otras cuatro energías de la Campaña 1-- queda como trabajo futuro; ver la sección de conclusiones
del informe técnico de la Campaña 1 para más detalle.

\subsection{Nombre de carpeta mal etiquetado}

Solo una de las dos carpetas de datos crudos tiene un nombre que no corresponde a la energía real de
la corrida (probablemente arrastrado de una etapa anterior de generación de los archivos de flujo de
entrada): \codigo{1keV/} (dentro de la Campaña 1) contiene en realidad la corrida de \textbf{10\,keV}
--verificado de forma independiente por dos vías (columna de energía cinética inicial en
\codigo{interaccion-completa-neutrones.txt} y el log crudo de Geant4; ver la sección
``Configuración de la simulación'' del informe técnico de la Campaña 1 para el detalle). La carpeta
\codigo{25meV/} de la Campaña 2 \textbf{sí} está correctamente nombrada (verificado: 25\,meV). Se
documenta aquí y en el informe técnico correspondiente para que no se propague el error;
\textbf{la carpeta \codigo{1keV/} no se renombró en disco} --los notebooks existentes leen rutas
relativas a ese nombre.

\subsection{Volumen activo del tanque (actualización del 18 de septiembre de 2026)}

Desde esta fecha, los 20 notebooks de la Campaña 1 subdividen la categoría ``Capturado'' según la
posición de la captura respecto a la caja del volumen activo del tanque
($z\in[0,1330]\,\text{mm}$, $|x|,|y|\leq480\,\text{mm}$, dada explícitamente -- no auto-detectada):
\textbf{Capturado (dentro del tanque)} y \textbf{Capturado fuera del tanque}. La fracción fuera del
tanque resultó marginal (0.03\,\%--0.4\,\% de los neutrones incidentes) en las 20 corridas, así que
esta separación no cambia ninguna cifra de captura total ya reportada -- solo añade precisión
metodológica. El informe técnico de la Campaña 1 documenta el detalle completo.

\section{Estructura del directorio \texttt{analysis/}}

\begin{verbatim}
analysis/
|-- notebooks/
|   `-- campana-2_con-S-alpha-beta_25meV/
|         Notebooks de la Campaña 2 (con S(alpha,beta), 25 meV)
|-- scripts/           Programas reproducibles
|-- wcd_analysis/       Funciones reutilizables (pendiente)
`-- tests/              Pruebas de las funciones de análisis (pendiente)
\end{verbatim}

Los notebooks de la Campaña 1 (cinco energías, sin $S(\alpha,\beta)$) no están replicados en este
repositorio: viven junto a sus datos crudos, cuyo tamaño (cientos de MB a varios GB por corrida) los
hace poco prácticos para versionar en git. El informe técnico de la Campaña 1 describe su convención
de nombres y su flujo de trabajo en la sección ``Reproducibilidad''.

Los programas definitivos deben aceptar rutas y parámetros por línea de comandos o archivos de
configuración; no deben contener rutas personales fijas.

\end{document}
